\section{Deligne's explicitly stated problem and the geometric weight boundary}
\label{sec:deligne-problem}

\subsection{D-WII-1.8.15: good degenerations of mixed Hodge structures}

There is an actual question next to the cited local-monodromy theorem.
Deligne (1980), \emph{La conjecture de Weil II}, 1.8.14--1.8.15,
printed pp.~178--179 (PDF pp.~42--43), asks to
``D\'egager une classe de `bonnes' variations de structures de Hodge mixtes''
on the punctured disc. In precise mathematical terms, for a good
unipotent variation $(V,W,F)$ with $N=\log T$, he asks for a finite
increasing filtration $M$ such that
\[
 NM_a\subseteq M_{a-2},\qquad
 N^b:\operatorname{gr}^{M}_{a+b}\operatorname{gr}^{W}_a V
       \xrightarrow{\sim}
       \operatorname{gr}^{M}_{a-b}\operatorname{gr}^{W}_a V
 \quad(b\geq0).
\]
He also asks for an asymptotic nilpotent orbit and for the filtration
$M$ to supply a mixed Hodge structure for the nilpotent orbit.
This is a faithful formal restatement of the displayed problem, not
an attribution of a question about the Riemann zeta function.

\paragraph{Answer: known geometric theorem, independent of the adopted RH.}
Admissible graded-polarizable variations provide the requested class.
For the precise definition use Kashiwara (1986), \S1.8,
conditions (1.8.1)--(1.8.4), pp.~993--994, and \S1.9: quasi-unipotence,
relative monodromy, and locally free extension of the Hodge and weight
graded bundles. Proposition 4.5.1 and Proposition 5.2.1,
pp.~1008--1010, identify the limiting data and the mixed Hodge
structure for $M$. His Theorem 4.4.1 and Theorem 4.5.2,
pp.~1007--1008, establish the compatibility of relative filtrations
and the normal-crossing extension criterion. These are theorems about
the defined class; the assertion is not that every variation satisfying
only Deligne's initial three axioms is automatically admissible.

For an explicit asymptotic estimate, Pearlstein (2000),
arXiv:math/0002030v3, Theorem 5.1(i)--(iii), p.~30,
gives a limiting filtration $F_\infty$ and the horizontal orbit
$F_{\rm nilp}(z)=\exp(zN)F_\infty$, with
\[
 d\bigl(F(z),F_{\rm nilp}(z)\bigr)
 \leq K(\operatorname{Im}z)^\beta
                  e^{-2\pi\operatorname{Im}z}
 \quad(\operatorname{Im}z\gg0)
\]
in the mixed Hodge metric. This supplies the asymptotic part in a
stronger quantitative form. The relative-monodromy mixed Hodge
structure and this estimate together answer the stated existence
problem. The question is settled by the cited Hodge-theoretic work;
no new proof or historical priority claim is made here.

\subsection{D-WII-1.9.5: the arithmetic base is a different gate}

At the end of 1.9.5, printed p.~181 (PDF p.~45), Deligne contrasts
normal-crossing divisors in a regular scheme of finite type over
$\mathbb Z$ with the preceding finite-field theory, already mentioning
$\operatorname{Spec}\mathbb Z[1/\ell]$. This observation must be
separated from the solved complex-analytic problem 1.8.15.
For the proper smooth local-field case a precise modern statement is
Scholze (2012), \emph{Perfectoid spaces}, Conjecture 9.3, p.~48
of arXiv:1111.4914:
for $H^i(X_{\overline K},\mathbb Q_\ell)$ the $j$th monodromy-graded
piece has Frobenius weight $i+j$. His Theorem 9.6, p.~49, proves
this in its stated set-theoretic complete-intersection situation
inside a projective smooth toric variety. We retain all of those
hypotheses; it is not a theorem for arbitrary proper smooth $X/K$.

BR-WM-01 identifies a concrete missing datum: a Frobenius
characteristic polynomial and the covariance equation do not determine
the monodromy filtration. Hence the adopted global zero-location
statements do not supply this datum through a reconstruction from
Euler polynomials. This is a proved failure of that proposed route,
not a metamathematical independence claim. The remaining gate for a
specified degeneration is the purity of each actual monodromy-graded
piece, with its actual $N$ and geometric realization.

\subsection{The universal stable curve: purity of a complex and weights of a fibre}

The RESIDUE intake, \S4(1), mixes two valid constructions with an
invalid identification. For a proper map from a smooth scheme over
a finite field, the constant complex is pure, and proper pushforward
preserves purity in the sense of mixed complexes. The exact source is
Beilinson, Bernstein, and Deligne (1982), Stabilit\'es 5.1.14 and its dual
form, Corollary 5.4.4, Theorem 5.4.5, and Corollary 5.4.6,
pp.~128--129 and 142. The splitting statements there hold after
base change to the algebraic closure: arithmetic Frobenius
semisimplicity is not silently added.

The universal stable-curve application uses Sun's stack decomposition
theorem, with its finite-diagonal and affine-stabilizer hypotheses,
as proved in CC-ST01 below. A splitting does not in general descend
merely from its existence on scheme charts. It does not say that every ordinary
cohomology group of a singular geometric fibre is pure of its
cohomological degree. BR-NODE-01 computes the actual group: its
weight-zero graph subspace and weight-one normalization part need
not have a single weight. Nearby smooth cohomology is a different
object. In the one-node genus-one smoothing its monodromy filtration
has graded pieces of weights zero and two, while the nodal special
fibre has only a one-dimensional weight-zero $H^1$.

For that last nearby-cycle statement the exact classical input is
SGA~7~II, Expos\'e XV, Th\'eor\`eme 3.4(ii)--(iii), printed
pp.~195--196: specialization, vanishing cycles, and the
Picard--Lefschetz inertia formula. The determinant
$\det H^1=\mathbb Q_\ell(-1)$ identifies the quotient twist once
the trivial specialization line has been identified. This explicit
case needs no general mixed-characteristic weight--monodromy theorem.
The entrance-path claims RESIDUE H.3, H.5, H.6 have their own
categorical proofs, unavailable in the supplied packet. Their
reported status is carried as an intake report, not promoted to a
newly checked theorem by this geometric calculation.
